\documentclass[11pt]{article}
\usepackage{amsmath,amssymb}
\usepackage[margin=1in]{geometry}
\usepackage{hyperref}
\emergencystretch=3em

\title{The phase-dressed normal moment}
\author{Claude, with Daniel Murfet}
\date{2026-09-07}

\begin{document}
\maketitle
\sloppy

\begin{abstract}
Unit 201 (programme~A, step~3). The one-dimensional normal moment dressed by a constant phase,
\[
S^{(\beta)}_{p/2}(a)=\int_0^\infty s^{p-1}e^{-\beta s^2+\beta as}\,ds
\]
(\texttt{phaseMoment β p a}; the paper's $S_{p/2}(a)$, kernel \texttt{quadKernel β a}), is
formalised with everything the assembly unit needs: integrability with the uniform Gaussian
domination $s^{p-1}e^{-\beta s^2+\beta as}\le e^{\beta R^2/2}s^{p-1}e^{-\beta s^2/2}$ for $|a|\le R$,
continuity in $a$ (dominated convergence on a ball), the Gaussian moments
$\int_0^\infty s^{p-1}s^je^{-\beta s^2}\,ds=\tfrac12\Gamma(\tfrac p2+\tfrac j2)\beta^{-(p/2+j/2)}$,
the zero-phase value $\tfrac12\Gamma(p/2)\beta^{-p/2}$, and the \emph{finite Taylor approximation}
\[
\Bigl|S^{(\beta)}_{p/2}(a)-\sum_{j<2K}\frac{(\beta a)^j}{j!}\cdot\tfrac12\Gamma(\tfrac p2+\tfrac j2)\beta^{-(p/2+j/2)}\Bigr|
\le\tau_K(\beta,b)\int_0^\infty s^{p-1}e^{-\beta s^2/4}\,ds\qquad(|\beta a|\le b),
\]
with the uniform tail constant $\tau_K$ of unit~199. Zero \texttt{sorry}/\texttt{axiom}; 9078 jobs.
\end{abstract}

\section{The things formalised}
{\small
\begin{verbatim}
noncomputable def phaseMoment (β p a : ℝ) : ℝ :=
    ∫ s in Ioi 0, s ^ (p - 1) * quadKernel β a s
noncomputable def gaussTail (β p : ℝ) : ℝ := ∫ s in Ioi 0, s ^ (p - 1) * exp (-(β/4 * s^2))
theorem gauss_integrableOn (hc : 0 < c) (hq : -1 < q) :
    IntegrableOn (fun s => s ^ q * exp (-(c * s^2))) (Ioi 0)
theorem quadKernel_le_of_abs_le (hβ) (hs : 0 ≤ s) (haR : |a| ≤ R) :
    quadKernel β a s ≤ exp (β R²/2) * exp (-(β/2 s²))
theorem phaseMoment_integrand_integrableOn (hβ) (hp : 0 < p) :
    IntegrableOn (fun s => s ^ (p - 1) * quadKernel β a s) (Ioi 0)
theorem continuous_phaseMoment (hβ) (hp) : Continuous fun a => phaseMoment β p a
theorem gaussMoment_eq (hβ) (hq : -1 < q) :
    ∫ s in Ioi 0, s ^ q * exp (-(β s²)) = Γ ((q+1)/2) * β ^ (-((q+1)/2)) / 2
theorem gaussMoment_nat_eq (j : ℕ) (hβ) (hp) :
    ∫ s in Ioi 0, s ^ (p-1) * (s ^ j * exp (-(β s²))) = Γ (p/2 + j/2) * β ^ (-(p/2 + j/2)) / 2
theorem phaseMoment_zero (hβ) (hp) : phaseMoment β p 0 = Γ (p/2) * β ^ (-(p/2)) / 2
theorem phaseMoment_taylor_le (hβ) (hp) (hab : |β * a| ≤ b) (K : ℕ) :
    |phaseMoment β p a - ∑ j ∈ range (2K), (β a)^j / j! * (Γ (p/2 + j/2) * β^(-(p/2+j/2)) / 2)|
      ≤ tailCoeff β b K * gaussTail β p
\end{verbatim}
}

\section{Mathematics}
The finite sum is rewritten as one integral, $\int_0^\infty s^{p-1}\bigl(\sum_{j<2K}(s\beta a)^j/j!\bigr)e^{-\beta s^2}\,ds$,
via the Gaussian moments (Mathlib's \texttt{integral\_rpow\_mul\_exp\_neg\_mul\_rpow} at $p=2$);
the difference with $S(a)$ is then a single integral whose integrand is bounded pointwise by
\texttt{phase\_tail\_le}, and \texttt{norm\_integral\_le\_of\_norm\_le} finishes. Nothing here
depends on the exponent data $(h,k)$; in the assembly the phase $a=\xi(\pi u)$ varies over the face
and the bound is uniform because $|\beta\xi|\le\beta\|\xi\|_\infty=:b$ on the cube.

\section{Lean notes}
\begin{itemize}
\item Dot notation \texttt{(h.const\_mul c).congr\_fun} fails with ``\texttt{And.congr\_fun}'':
  \texttt{Integrable} unfolds to a conjunction, so the result of \texttt{Integrable.const\_mul}
  is not seen as an \texttt{IntegrableOn}. Ascribe the type in a \texttt{have} first.
\item \texttt{integral\_finset\_sum}/\texttt{integrable\_finset\_sum} are deprecated for
  \texttt{integral\_finsetSum}/\texttt{integrable\_finsetSum}.
\item \texttt{IntegrableOn.congr\_fun} hands the pointwise goal as an applied lambda;
  \texttt{beta\_reduce} before \texttt{rw}.
\item \texttt{continuousAt\_of\_dominated} with the bound depending on the centre $a_0$ (radius-one
  ball, $R=|a_0|+1$) is the clean way to get continuity of a parametric integral whose integrand is
  not globally dominated.
\end{itemize}

\section{Next}
Unit 202: the leading-term theorem. Define
$C=\frac{1}{(m-1)!\prod_{i\in J}k_i}\int_{(0,1]^d}\eta(\pi u)\,S^{(\beta)}_{p/2}(\xi(\pi u))\,w(u)\,du$
and prove $I_N/(N^{-p}(\log N)^{m-1})\to C$ by \texttt{tendsto\_of\_approx}: the finite Taylor
approximants converge by unit~200, the $N$-side error is bounded by the bare monomial theorem at
$\beta/4$ times $\tau_K$, and $|C-P_K|$ by \texttt{phaseMoment\_taylor\_le}.
\end{document}
